\chapter{Introduction}\label{sec:introduction}

\section{The Sherrington--Kirkpatrick model and the Parisi formula}\label{sec:intro-background}

Let $(g_{ij})_{1\le i<j\le N}$ be independent standard Gaussian
variables. The Sherrington--Kirkpatrick (SK)
model~\cite{SherringtonKirkpatrick1975} at zero external field and
inverse temperature $\beta>0$ has Hamiltonian
\[
 H_N(\sigma)=\frac1{\sqrt N}\sum_{1\le i<j\le N}g_{ij}\sigma_i\sigma_j,
 \qquad \sigma\in\{-1,1\}^N.
\]
Its logarithmic partition function and expected pressure are
\[
 F_{N,\beta}=\log\sum_{\sigma\in\{-1,1\}^N}e^{\beta H_N(\sigma)},
 \qquad
 p_{N,\beta}=\frac{\E F_{N,\beta}}N.
\]
For a probability measure $\mu$ on $[0,1]$ with distribution function
$\alpha(q)=\mu([0,q])$, let $u_\alpha$ be the solution of the Parisi
equation on $[0,1]\times\R$, and define the Parisi functional
$\mathcal P_\beta$ and the Parisi value $P_\beta$ by
\begin{equation}\label{eq:intro-parisi}
\begin{gathered}
 \partial_qu_\alpha=-\frac{\beta^2}{2}
 \bigl(\partial_{xx}u_\alpha+\alpha(q)(\partial_xu_\alpha)^2\bigr),
 \qquad u_\alpha(1,x)=\log(2\cosh x),\\
 \mathcal P_\beta(\alpha)=u_\alpha(0,0)
       -\frac{\beta^2}{2}\int_0^1q\alpha(q)\dd q,
 \qquad P_\beta=\inf_{\mu\in\mathcal P([0,1])}\mathcal P_\beta(\alpha),
\end{gathered}
\end{equation}
where $\mathcal P([0,1])$ is the set of probability measures on
$[0,1]$. For a general $\mu$, the equation is understood in the weak
sense of~\cite[Theorem~4]{JagannathTobasco2016}.

Sherrington and Kirkpatrick introduced this model and studied it by a
replica-symmetric calculation, which fails at low
temperature~\cite{SherringtonKirkpatrick1975}; a longer treatment, with
numerical comparisons, is given in~\cite{KirkpatrickSherrington1978}.
De~Almeida and Thouless showed that the replica-symmetric saddle point
becomes unstable~\cite{deAlmeidaThouless1978}. Parisi proposed a
hierarchical breaking of the replica symmetry in~\cite{Parisi1979}, and
he developed a sequence of approximate solutions and the order
parameter on $[0,1]$ that leads to \eqref{eq:intro-parisi}
in~\cite{Parisi1980Sequence,Parisi1980OrderParameter}. He later
interpreted this order parameter through the distribution of overlaps
between states~\cite{Parisi1983OrderParameter}. M\'ezard, Parisi,
Sourlas, Toulouse, and Virasoro then described the organization of the
states, with sample-dependent weights and ultrametric overlaps,
in~\cite{MezardEtAl1984Nature}.

The Parisi formula states that $p_{N,\beta}\to P_\beta$ as
$N\to\infty$. Guerra proved the upper bound $p_{N,\beta}\le P_\beta$ for
every $N$~\cite{Guerra2003}, and Talagrand proved the matching lower
bound in the limit~\cite{Talagrand}. The infimum in
\eqref{eq:intro-parisi} is attained at a unique measure
$\mu_\beta$~\cite{AuffingerChen2015Uniqueness}, the Parisi measure,
whose distribution function we denote by $\alpha_\beta$. For $\beta>1$,
its support is an interval $[0,q_\beta]$ with
$0<q_\beta<1$~\cite[Theorem~1.1]{Lopatto}; for $\beta>1$ close to $1$,
this was proved earlier by Zhou~\cite{Zhou2026FRSB}.

\section{Finite-size corrections}\label{sec:intro-finite-size}

In this paper, we study the rate of convergence in the Parisi formula at
low temperature. By Guerra's bound, the \emph{mean bias}
$P_\beta-p_{N,\beta}$ is nonnegative for every $N$. Up to a factor
$1/\beta$, it is the finite-size correction to the expected free energy
per spin (\cref{sec:intro-results}).

At high temperature, the mean bias is small. For $0<\beta<1$, Guerra
and Toninelli showed by a quadratic coupling of two replicas that it is
of order at most $N^{-1}$~\cite[Section~3]{GuerraToninelli2002Quadratic}.
With a nonzero external field, Lopatto proved that the limiting
pressure equals the replica-symmetric value up to and including the
de~Almeida--Thouless line~\cite[Theorem~1.1]{Lopatto2026ATLine}, and
Kusuoka and Nakajima proved that the pressure is within $O(N^{-1})$ of
this value on compact subsets of the strict de~Almeida--Thouless
region~\cite[Theorem~1.1]{KusuokaNakajima2026}.

Below the critical temperature, the mean bias is expected to decay more
slowly. In~\cite[Theorem~1.1, (1.5)]{AL}, we proved that
$0\le P_\beta-p_{N,\beta}\le C\beta N^{-1/2}$ for every $\beta>0$ and
$N\ge1$, with a numerical constant $C$. For fixed $\beta>1$, we used
the interval support of $\mu_\beta$ to prove the sharper upper bound
$C_\beta N^{-2/3}$ and the lower bound $c_\beta N^{-11/12}$, for all
sufficiently large $N$~\cite[Theorem~1.4]{AL}. In physics, the
correction is expected to decay like $N^{-\omega}$ with $\omega<1$ at
low temperature. Two values have been proposed for the full model,
$\omega=2/3$ and $\omega=3/4$. A calculation retaining only
longitudinal fluctuations gives
$\omega=4/5$~\cite{ParisiSarraTalamanca2019}. We review these predictions in
\cref{subsec:intro-physics}. We improve the exponent of the upper bound
from $2/3$ to $7/10$.

\section{Main result}\label{sec:intro-results}

Our main result is the following.

\begin{theorem}\label{thm:main}
For every $\beta>1$ there is a constant $C_\beta<\infty$ such that
\[
 P_\beta-p_{N,\beta}\le C_\beta N^{-7/10}
 \qquad(N\ge2).
\]
\end{theorem}

Since $P_\beta-p_{N,\beta}\ge0$ for every $N$~\cite{Guerra2003},
\cref{thm:main} is a two-sided estimate. In physical units, the
expected free energy density is $\E[-F_{N,\beta}/(\beta N)]$, and its
limit is $-P_\beta/\beta$. By \cref{thm:main}, we have
\[
 0\le \E\Bigl[-\frac{F_{N,\beta}}{\beta N}\Bigr]
     -\Bigl(-\frac{P_\beta}{\beta}\Bigr)
 =\frac{P_\beta-p_{N,\beta}}{\beta}
 \le\frac{C_\beta}{\beta}N^{-7/10}
 \qquad(N\ge2).
\]
For the expected logarithm of the partition function itself, the
corresponding bound is $0\le NP_\beta-\E F_{N,\beta}\le C_\beta N^{3/10}$. The constant
$C_\beta$ depends on $\beta$, and we do not control it as
$\beta\downarrow1$ or as $\beta\to\infty$. In particular, we do not
obtain a corresponding bound for the ground-state energy.

Together with the lower bound in~\cite[Theorem~1.4]{AL},
\cref{thm:main} implies that for every fixed $\beta>1$,
\[
 c_\beta N^{-11/12}\le P_\beta-p_{N,\beta}\le C_\beta N^{-7/10}
\]
for all sufficiently large $N$. The exponent $7/10$ lies between the
exponent $2/3$ of~\cite[Theorem~1.4]{AL} and the exponent $3/4$
predicted by Franz, Parisi, and
Ricci-Tersenghi~\cite{FranzParisiRicciTersenghi2026}, which we discuss
in \cref{subsec:intro-physics}. In \cref{subsec:intro-openai}, we
compare \cref{thm:main} with two conditional consequences of results
stated in the OpenAI
manuscripts~\cite{OpenAIFluctuationScale,OpenAILimitingLaw}. An
uncentered comparison stated there would give an unspecified strict
improvement over $2/3$, and the variance asymptotic stated there would
imply the lower bound $P_\beta-p_{N,\beta}\ge c_\beta N^{-5/6}$.

We now outline the proof; \cref{sec:overview} gives a detailed
overview. We follow the interpolation scheme of~\cite{AL}, which is a
quantitative form of the Guerra interpolation~\cite{Guerra2003} that
Talagrand used for the positive fractional moments of the partition
function~\cite{Talagrand2007FractionalMoments}. For a tilt
$0<\lambda\le\lambda_0$, let $e_{N,\beta}(\lambda)\ge0$ be the error,
in energy units, between the constrained Parisi value and the normalized
logarithmic moment of order $\lambda/\beta$ of the completed partition
function. The constraint requires mass at least $\lambda/\beta$ at zero
(\cref{sec:fnd}).
Imposing this mass raises the Parisi value by at most
$C_\beta\lambda^5$ in energy units~\cite[Lemma~7.2]{AL}. Using this
bound and the convexity of Gaussian logarithmic
moments~\cite[Theorem~1]{ChenConvexity}, as in the proof
of~\cite[Theorem~7.8, Step~1]{AL}, we find
\[
 0\le\frac{P_\beta-p_{N,\beta}}{\beta}
 \le C_\beta\lambda^5+2e_{N,\beta}(\lambda).
\]
For sufficiently large $N$, we take $\lambda=N^{-7/50}$, so it
suffices to show that
$e_{N,\beta}(N^{-7/50})\le C_\beta N^{-7/10}$.

Up to a factor of four, $e_{N,\beta}(\lambda)$ is the integral over
$\rho\in(0,1]$ of a defect $\calD_N(\rho)=\beta(A+2J+V)$ along the
interpolation, evaluated at a minimizing profile of an $N$-dimensional
variational problem. The terms $A$ and $J$ are controlled directly
(\cref{sec:fnd,sec:th}). The difficulty is $V$, the integrated variance
of a self-overlap. For $\rho$ above a cutoff $\rho_0$, we bound $V$ by
a one-site cavity argument (\cref{part:cavity}). Removing one spin leads
to an inequality $V\le(1-\rho)(1+\eps_N)V+{}$errors, where
$1+\eps_N$ bounds a scalar response operator at the minimizing profile,
and this inequality contracts once $\eps_N\le\rho/4$. To prove this
bound, we show that the response operator at the Parisi measure is
bounded by the identity, which reflects the marginal stability of
$\mu_\beta$, and that the minimizing profile is close to a scalar trial
profile (\cref{part:profile}). For $\rho\le\rho_0$, we bound $V$ by
the quantile-shell estimates of~\cite[Section~3]{AL}, refined in
\cref{sec:wi}, and by a localization of the lowest quantiles of the
minimizer (\cref{sec:lz-root}), which refines~\cite[Lemma~7.6]{AL}. The cutoff is
$\rho_0=\min\{1,K_\beta\eta^{1/7}\}$, where the trial budget
$\eta=e_{N,\beta}(\lambda)+C_\beta\lambda^5$ contains the unknown error.
The result is an implicit inequality of the form
\[
 e_{N,\beta}(\lambda)\lesssim N^{-2/3}
 \bigl(e_{N,\beta}(\lambda)+\lambda^5\bigr)^{1/21}+\text{smaller terms}.
\]
If $e_{N,\beta}(\lambda)$ and $\lambda^5$ are both of order $N^{-a}$,
this inequality is balanced when $a=2/3+a/21$, that is, $a=7/10$.

\section{Relation to physical predictions}\label{subsec:intro-physics}

\paragraph{Expectations and early calculations.}
Below the critical temperature, physicists expect the finite-size
correction to the free energy density to be nonanalytic, of order
$N^{-\omega}$ with $\omega<1$. This expectation is related to marginal
stability. In a truncated replica model near criticality,
De~Dominicis and Kondor found that the stability matrix of the Parisi
solution has nonnegative eigenvalues and many zero
modes~\cite[pp.~606--608]{DeDominicisKondor1983}. These zero modes make the
expansion in fluctuations around this solution singular;
see~\cite[Section~2]{FranzParisiRicciTersenghi2026}. Parisi, Ritort,
and Slanina studied finite-size corrections at
criticality~\cite{ParisiRitortSlanina1993Critical},
where the free energy density has a correction $(\log N)/(12N)$
(see~\cite{Billoire2006FiniteSize}), as well as slightly below
criticality and along the de~Almeida--Thouless line, including
corrections to overlap distributions and to
ultrametricity~\cite{ParisiRitortSlanina1993Corrections}.

\paragraph{The exponent $2/3$.}
At zero temperature, Bouchaud, Krzakala, and Martin found numerically,
for Gaussian couplings, that the mean ground-state energy per spin
approaches its limit like
$N^{-2/3}$~\cite[Section~VI\,B]{BouchaudKrzakalaMartin2003}. In a
study of the SK model with binary couplings, Boettcher fits the corrections to the
mean of the ground-state energy per spin and to the width of its
distribution with exponents close to $2/3$ and $3/4$,
respectively~\cite[Section~3]{Boettcher2005}. In a later study of
diluted versions of the model, the fit for the undiluted model is
again close to $2/3$~\cite[Table~I]{Boettcher2020Diluted}. At positive
temperature, Billoire estimated the correction numerically for binary
couplings, by integrating overlap observables along the interpolation of
Guerra and Toninelli~\cite{GuerraToninelli2002}, and found data
compatible with $\omega=2/3$~\cite{Billoire2006FiniteSize}. Aspelmeier,
Billoire, Marinari, and Moore derived $\omega=2/3$ from a heuristic
argument: a system of $N$ spins resolves about $N^{1/6}$ steps of
replica symmetry breaking, and truncation after $K$ steps changes the
free energy density by order
$K^{-4}$~\cite[Section~2]{AspelmeierEtAl2008FiniteSize}.

\paragraph{Longitudinal fluctuations and the exponent $4/5$.}
Parisi and Biscari obtained an earlier contribution of order
$N^{-3/4}\log N$ from the positive longitudinal eigenvalues near
criticality, neglecting interactions between
modes~\cite[Section~1 and Eq.~(5.25)]{ParisiBiscari1992}. Their
contribution from the longitudinal zero modes diverges when treated separately
from the omitted transverse modes~\cite[discussion following
Eq.~(5.27)]{ParisiBiscari1992}; see
\cite[Section~1]{ParisiSarraTalamanca2019}. Parisi, Sarra,
and Talamanca later used a resummation-based scaling argument for the
longitudinal fluctuations of a truncated free-energy functional for the SK model near
criticality and predicted
$\omega=4/5$~\cite[Section~3, Eqs.~(16)--(19)]{ParisiSarraTalamanca2019}.
They omit the transverse fluctuations, so this exponent describes their
longitudinal approximation. To determine the exponent for the full model, one
must also control the omitted modes.

\paragraph{The exponent $3/4$.}
Franz, Parisi, and Ricci-Tersenghi~\cite{FranzParisiRicciTersenghi2026}
argue that $\omega=2/3$ is incorrect and that the correct value is
$\omega=3/4$, possibly up to logarithmic
factors~\cite[Eq.~(14)]{FranzParisiRicciTersenghi2026}. Their
calculation for the Gaussian SK model matches two regimes: one in which the
overlaps of two coupled systems lock and the Guerra--Toninelli
integrand decays like $(Nt)^{-1/2}$, and a perturbative regime in which
this integrand is of order $t^{-3/2}/N$~\cite[Eqs.~(9)--(12) and
(15)--(17)]{FranzParisiRicciTersenghi2026}. They estimate the large-$z$ locking
exponent semianalytically by sampling random ultrametric
trees generated by the Bolthausen--Sznitman coalescent and extrapolating
in the number of leaves~\cite[Materials and Methods]{FranzParisiRicciTersenghi2026}.
Matching this behavior to the perturbative regime gives the scaling
prediction. The
fit of the energy uses random regular graphs of degree
four~\cite[Eq.~(18) and Fig.~2]{FranzParisiRicciTersenghi2026}, and the
simulations of the mode-locking integral use rewired random regular
graphs of degree four with binary couplings (see the subsection
``Direct simulations of the mode-locking integral''
of~\cite[Materials and Methods]{FranzParisiRicciTersenghi2026}). These
models differ from the Gaussian SK model in their connectivity, and the
latter also in the law of the couplings.

\paragraph{Consequences of \cref{thm:main}.}
\Cref{thm:main} implies that $P_\beta-p_{N,\beta}=O(N^{-7/10})$ for
every fixed $\beta>1$; in particular, the correction is $o(N^{-2/3})$. We therefore exclude a nonzero leading
correction of order $N^{-2/3}$ for the Gaussian SK model at every fixed
temperature below the critical one. This is
consistent with the prediction of Franz, Parisi, and Ricci-Tersenghi.
\Cref{thm:main} does not determine the leading exponent $\omega$. Since our
constants are not uniform as $\beta\to\infty$, it also says nothing
about the numerical results at zero temperature described above.

The estimates on which our proof rests have a similar structure with
two regimes, although they concern a different interpolation. If we
ignore logarithms and the dependence on the tilt $\lambda$, the defect
$\calD_N(\rho)$ is at most of order $(N\rho)^{-2/3}$ for small $\rho$
and of order $(N\rho^2)^{-1}$ on the interval where our cavity argument
applies. If the second bound held for all $\rho\ge N^{-1/4}$, where the
two bounds are equal, integration would give $N^{-3/4}$. In our proof,
the cavity interval begins only at $\rho_0\ge N^{-1/10}$, and the tilt
$\lambda=N^{-7/50}$ enters both bounds; these two features account for
the difference between $7/10$ and $3/4$.

Marginal stability also enters our proof, through a scalar statement:
at the Parisi measure, the multiplier that dominates our scalar response
operator equals one on $\supp\mu_\beta$ (\cref{sec:sc}). Crisanti and
Rizzo wrote the corresponding identity in physical variables, on the
part of the order parameter that
varies~\cite[Eqs.~(22)--(23)]{CrisantiRizzo2002}. This scalar statement
does not describe the full stability spectrum of the replica
calculation.

\paragraph{Mean, width, and large deviations.}
The exponents $\omega$ for the finite-size correction discussed above concern
the mean. The sample-to-sample fluctuations
of the free energy and its large deviations are governed by different
exponents. In a truncated replica model near criticality, Kondor
obtained a fifth-order dependence on $n$ along the Parisi branch for small positive
$n$~\cite[pp.~L129--L130, Eq.~(11)]{Kondor1983}. Interpreted
as a formal prediction for the SK model, this says that the
replica-normalized pressure in the thermodynamic limit,
$\lim_{N\to\infty}(Nn)^{-1}\log\E e^{nF_{N,\beta}}$, differs from
$P_\beta$ first at order $n^5$ as $n\downarrow0$ with
$n>0$ (see~\cite[pp.~1 and~3]{ParisiRizzo2008}).
Parisi and Rizzo extended this calculation throughout the
low-temperature phase in~\cite{ParisiRizzo2008}
and~\cite[Appendix~B]{ParisiRizzo2009}. Equivalently, in this formal
expansion, the limit of $N^{-1}\log\E e^{n(F_{N,\beta}-NP_\beta)}$
begins at order $n^6$, and its Legendre transform gives a rate with
power $6/5$ for upward deviations of
$F_{N,\beta}/N$~\cite{ParisiRizzo2008}.
Crisanti, Paladin, Sommers, and
Vulpiani interpreted the replica number
as a tilt over the disorder and predicted that the free energy per spin
fluctuates on the scale
$N^{-5/6}$~\cite{CrisantiEtAl1992Fluctuations}. Under an asymmetric
scaled-law assumption, the same calculation also predicts a correction to the mean
density on this scale~\cite[p.~1331, Eq.~(34) and following
text]{CrisantiEtAl1992Fluctuations}. This conclusion uses the
truncated replica theory near criticality and a matching of the disorder
tails to the typical peak.
In our convention, in which $F_{N,\beta}$ is the logarithm of the partition
function, Parisi and Rizzo's
replica predictions give upward density deviations with speed $N$ and
power $6/5$, and downward density deviations with speed $N^2$ and power
$12/5$~\cite[Eqs.~(12)--(15)]{ParisiRizzo2010}. Smoothly matching these
tails to the finite-size peak predicts the typical density scale
$N^{-5/6}$. This matching is an additional assumption; the limits
$n\downarrow0$ and $N\to\infty$ cannot simply be
interchanged~\cite[Section~VI]{ParisiRizzo2009}.
Aspelmeier, Billoire,
Marinari, and Moore also predicted a variance of order $N^{1/3}$ for
$F_{N,\beta}$~\cite{AspelmeierEtAl2008FiniteSize}.
Aspelmeier derived exact Gaussian interpolation identities relating
the variance of the free energy to two- and four-replica overlaps between
systems with correlated disorder~\cite[Eqs.~(14) and~(31)]{Aspelmeier2008Exact}.
Combining these identities with replica estimates of the bond-chaos
overlap distribution gives a predicted upper bound $1/4$ for the
extensive standard-deviation exponent~\cite[Section~4.3]{Aspelmeier2008Exact,Aspelmeier2008Bond}.
The identities are exact; this exponent estimate uses replica
approximations and concerns the width.
At zero temperature,
the width of the distribution of the ground-state energy has been
studied numerically in~\cite{Boettcher2005,Palassini2008,Boettcher2010Fluctuations}.
For Gaussian couplings, Palassini finds data consistent with a mean
correction exponent $2/3$ and, after subleading corrections, favors the
exponent $5/6$ over
$3/4$ for the standard deviation per spin~\cite[Sections~3--4 and~6]{Palassini2008}.
Boettcher discusses the difficulty of distinguishing the exponents
$3/4$ and $5/6$ after corrections to
scaling~\cite[Section~IV]{Boettcher2010Fluctuations}.
Monthus and Garel studied rare-disorder scenarios and the matching
between typical and large-deviation tails for the ground-state energy of the SK model at zero
temperature~\cite[Section~III C--F]{MonthusGarel2010}.
Their alternative width--tail relations illustrate the additional
assumptions involved in such matching.

\section{Related work}\label{sec:intro-related}

\paragraph{The Parisi formula.}
Besides the proofs of Guerra and Talagrand, the rigorous theory of the
Parisi formula includes the following results. Guerra and Toninelli
proved by interpolation that the thermodynamic limit
exists~\cite{GuerraToninelli2002}, and Aizenman, Sims, and Starr
formulated an extended variational principle in terms of random overlap
structures~\cite{AizenmanSimsStarr2003}. Panchenko proved that the
Ghirlanda--Guerra identities~\cite{GhirlandaGuerra1998} imply ultrametricity of
overlaps~\cite[Theorem~1]{Panchenko2013Ultrametricity}, and he proved
the Parisi formula for mixed $p$-spin models, including mixtures with
odd interactions~\cite{Panchenko2014Parisi}. We refer to the books of
Talagrand~\cite{TalagrandVolI2011,TalagrandVolII2011} and
Panchenko~\cite{Panchenko2013Book} for accounts of this theory;
Panchenko develops it through Ruelle probability cascades.

\paragraph{The Parisi measure.}
Toninelli proved that the replica-symmetric value is not optimal for
$\beta>1$ at zero field~\cite{Toninelli2002}. Auffinger and Chen proved
that the minimizer is unique~\cite{AuffingerChen2015Uniqueness}, that
zero belongs to its support, and that the largest point of its support
is an atom for $1<\beta\le3/2$~\cite[Theorems~1
and~4]{AuffingerChen2015Properties}. They also proved that its
distribution function is smooth on $[a,b)$ whenever $(a,b)$ is
contained in its support, for $0\le a<b\le1$~\cite[Theorem~2]{AuffingerChen2015Properties}.
They expressed the Parisi functional through a stochastic control
problem, which they used in their uniqueness proof, and they proved
martingale identities for the derivatives of
$u_\alpha$~\cite[Theorem~3 and Lemma~2]{AuffingerChen2015Uniqueness}. An earlier
Hamilton--Jacobi--Bellman representation appears
in~\cite[Section~6.2]{BovierKlimovsky2009}, and Jagannath and Tobasco
developed a dynamic programming treatment of the
functional~\cite[Lemma~18]{JagannathTobasco2016}. First-variation
formulas along segments and optimality conditions on the support are
given by Jagannath and Tobasco~\cite[Lemma~3.2 and
Proposition~1.1]{JagannathTobasco2017PhaseDiagram} and by
Chen~\cite[Theorem~2 and Proposition~1]{Chen2017Variational}. Lopatto
proved, in addition to the interval support for every $\beta>1$, that
$\mu_\beta$ is the sum of a density that is smooth up to the endpoints
of $[0,q_\beta]$ and an atom at $q_\beta$~\cite[Theorem~1.1]{Lopatto}.
Auffinger, Chen, and Zeng proved that the Parisi
minimizer at zero temperature has infinitely many support
points~\cite{AuffingerChenZeng2020}. H.-B.~Chen proved that
$q_\beta\to1$ as $\beta\to\infty$, and that the Stieltjes measure
associated with the Parisi order parameter at zero temperature has a
smooth density and support
$[0,1)$~\cite[Theorems~1.2 and~1.3]{Chen2026ZeroTemperatureSupport}. The
full support at zero temperature is also stated
in~\cite[Theorem~1.1]{OpenAIFullSupport}. We use the interval support
at fixed $\beta>1$ throughout, and the regularity of $\mu_\beta$ in
\cref{part:profile}.

\paragraph{Fluctuations of the free energy.}
The fluctuations of $F_{N,\beta}$ depend on $\beta$ and on the external
field. For $0<\beta<1$, Aizenman, Lebowitz, and Ruelle proved a
Gaussian limit for its fluctuations, which are of order
one~\cite{AizenmanLebowitzRuelle1987}, and Comets and Neveu studied this
limit by stochastic calculus~\cite{CometsNeveu1995}. With a nonzero
external field, Chen, Dey, and Panchenko proved that the variance is of
order $N$ and that a central limit theorem holds, under assumptions on
the interactions that include the SK
model~\cite[Theorems~1 and~2]{ChenDeyPanchenko2017}. At $\beta=1$, Chen
and Lam proved a variance bound of order
$\log^2N$~\cite[Theorem~1.1]{ChenLam2019}. Schertzer proved that the
variance is at most $(\log N)/4+C$ and at least of order
$\log\log\log N$~\cite[Theorems~1 and~2]{Schertzer2026Critical}, and Du
and Huang proved that $\Var(F_{N,1})=(\log N)/6+O(1)$ and that the
fluctuations have a Gaussian limit~\cite[Theorem~1.3]{DuHuang2026}. Dey and Kang obtained
variance asymptotics and a Gaussian limit as $\beta$ approaches $1$ from
below along sequences $\beta=\beta_N$ satisfying
$N^{1/3}(1-\beta_N^2)\to c\in(0,\infty]$~\cite[Theorem~1.1]{DeyKang2026}.
For fixed $\beta>1$, Chatterjee proved the superconcentration bound
$\Var(F_{N,\beta})\le C_\beta N/\log N$ and related it to disorder
chaos~\cite[Theorem~1.5]{Chatterjee2009Chaos}. In~\cite[Theorem~1.4]{AL},
we proved that $\Var(F_{N,\beta})$ lies between constant multiples of
$N^{4/15}$ and $N^{7/15}$ for large $N$. The OpenAI manuscripts state
the order $N^{1/3+o(1)}$, a positive limit for
$N^{-1/3}\Var(F_{N,\beta})$, and a limiting law for the normalized
fluctuations; see \cref{subsec:intro-openai}. A variance bound alone
does not bound the mean bias from above, but a lower bound on the
variance gives a lower bound on the mean bias
through~\cite[Lemma~7.4]{AL}. For the spherical SK model, Baik and Lee
proved that the free energy per spin has Gaussian fluctuations of order
$N^{-1}$ at high temperature and Tracy--Widom fluctuations of order
$N^{-2/3}$ at low temperature, governed by the largest eigenvalue of
the coupling matrix~\cite[Theorems~2.10 and~2.11]{BaikLee2016}.

\paragraph{Gaussian-process comparisons.}
Chu and Raginsky use Gaussian-process geometry and information-theoretic
tensorization to bound expected soft maxima~\cite{ChuRaginsky2026}.
Their application to the SK model gives a finite-volume comparison with the Parisi
variational value up to universal multiplicative
constants~\cite[Section~6, Theorem~10]{ChuRaginsky2026}.
It uses a uniform spin prior and normalizes the Hamiltonian differently.
Its conclusion is a multiplicative comparison; the present theorem
controls the additive error to the limiting value.

\paragraph{Fractional moments.}
Talagrand computed the limits of normalized logarithms of positive
moments of the partition function~\cite{Talagrand2007FractionalMoments}. For moments of order
$0<s<1$, the limit of $(Ns)^{-1}\log\E e^{sF_{N,\beta}}$ is the minimum
$P_\beta(s)$ of $\mathcal P_\beta$ over measures $\mu$ with
$\mu(\{0\})\ge s$. For the SK model at zero field, this formula is
contained in the one-level specialization of the multiscale formula of
Contucci and Mingione~\cite[Theorem~2.1]{ContucciMingione2019}. Chen,
Guionnet, Ko, Lacroix-A-Chez-Toine, and Mourrat stated it in this form
and used it to study one-sided large deviations of the ground-state
energy~\cite[Theorem~2.1]{ChenEtAl2026LargeDeviations}.
In~\cite[Lemma~7.2 and Proposition~7.3]{AL}, we proved that for fixed
$\beta>1$, the increase $P_\beta(s)-P_\beta$ is of order $s^5$ as
$s\downarrow0$. This is a rigorous counterpart of the replica
calculations of Kondor and of Parisi and Rizzo described in
\cref{subsec:intro-physics}, and Pei identified the leading coefficient
of this increase~\cite[Theorem~1(a)]{Pei2026}. The tilt
$\lambda=N^{-7/50}$ in our proof balances this cost against the
fractional-moment error.

\paragraph{Origins of the methods.}
The interpolation of~\cite[Section~2]{AL} is a finite-dimensional,
quantitative form of Guerra's interpolation, with which he proved the
upper bound in the Parisi formula~\cite{Guerra2003}. For posterior
estimates under its tilted laws, we combine the Brascamp--Lieb
inequality~\cite[Theorem~4.1]{BrascampLieb1976}, in its inverse-Hessian
form~\cite[Proposition~2.2]{AL}, with a joint-source argument from the
proof of~\cite[Lemma~2.3]{OpenAIFluctuationScale} (\cref{sec:gauss}).
Near the singular end of the interpolation, our estimates refine the
quantile-window bound of~\cite[Lemma~3.3]{AL} and the localization of
lowest quantiles in~\cite[Lemma~7.6]{AL}. We close the final implicit
inequality as in the proof of~\cite[Proposition~7.7]{AL}.

The cavity method originates in physics. An early precursor is the
single-site description of Thouless, Anderson, and
Palmer~\cite[Section~3, Eqs.~(10)--(12)]{ThoulessAndersonPalmer1977},
and M\'ezard, Parisi, and
Virasoro formulated the cavity construction for broken replica symmetry
in~\cite{MezardParisiVirasoro1986}; see also the exposition
in~\cite{MezardParisiVirasoro1987}. We use Talagrand's rigorous version
of the one-site argument~\cite[Section~1.6]{TalagrandVolI2011}. The
expansion under hierarchical Gaussian laws relies on normalized replica
trees and covariance insertions of the type developed
in~\cite[Section~2]{OpenAIFluctuationScale}. In the scalar stability
argument, the first variation of $\mathcal P_\beta$ under transports of
levels is due to Talagrand~\cite[Lemma~3.7]{Talagrand2006ParisiMeasures};
see~\cite[Section~3, (19)]{AuffingerChen2015Properties}. We compute the
second variation and compare it with a response multiplier, using
convexity of $\mathcal P_\beta$ along segments of distribution
functions~\cite{AuffingerChen2015Uniqueness}. This multiplier equals one
on $\supp\mu_\beta$, by an identity of
Talagrand~\cite{Talagrand2006ParisiMeasures}, in the form stated
in~\cite[Theorem~5]{AuffingerChen2015Properties}
and~\cite[Lemma~2.6]{Lopatto}, and by the interval
support~\cite[Theorem~1.1]{Lopatto}. Above the support, we bound it by one
using the curvature of the forward density and its heat evolution
from~\cite[Proposition~4.1 and Lemma~2.4(ii)]{Lopatto}.

\paragraph{Our related work.}
This paper continues~\cite{AL}, where we proved the quantitative Parisi
formulas recalled in \cref{sec:intro-finite-size}, the variance bounds
above, and upper-tail estimates for $F_{N,\beta}$ with exponent
$6/5$~\cite[Theorem~1.5]{AL}. The interval-support theorem used in both
papers was proved by the second author~\cite{Lopatto}.
In~\cite{AronowLopatto2026EA}, we studied a different approximation of
the SK pressure. On $d$-dimensional tori of side length $L\ge3$ and on
hypercubes, with independent Gaussian couplings of variance $1/(2d)$
and $1/d$, respectively, the expected pressure per site converges to the
SK pressure as $d\to\infty$, for every fixed inverse temperature
$\beta\ge0$ and every uniform field $h\in\R$; on tori, the convergence
is uniform in the side length. The present paper concerns the
finite-size bias of the SK model itself.

\section{Comparison with the OpenAI preprints}\label{subsec:intro-openai}

We obtained the exponent $7/10$ on September~23, 2026, following a
suggestion from Silvio Franz on September~15, 2026. The OpenAI
mathematics repository~\cite{OpenAIMathRepository} contains
\emph{The low-temperature Sherrington--Kirkpatrick fluctuation
scale}~\cite{OpenAIFluctuationScale} and its companion \emph{The
low-temperature Sherrington--Kirkpatrick free-energy limiting
law}~\cite{OpenAILimitingLaw}, both dated September~24, 2026. We refer
to their initial repository version.

\paragraph{Stated results.}
For every fixed $\beta>1$, \cite[Theorem~1.1]{OpenAIFluctuationScale}
states that $\Var(F_{N,\beta})=N^{1/3+o(1)}$. The companion paper states
that $N^{-1/3}\Var(F_{N,\beta})$ converges to a positive limit and that
$F_{N,\beta}$, standardized by its exact mean and variance, converges in
law to a nondegenerate limit~\cite[Theorems~1.1
and~1.2]{OpenAILimitingLaw}. These results concern fluctuations about
the finite-size mean $\E F_{N,\beta}$, and they do not by themselves
estimate $P_\beta-p_{N,\beta}$. We do not use these results, or the
comparison discussed next, in the proof of \cref{thm:main}.

\paragraph{A conditional upper bound.}
An improvement of the exponent $2/3$ for the mean bias is implicit in an
uncentered comparison stated in~\cite{OpenAILimitingLaw}. We use the
notation of that paper. Let $T=\beta^2$ and $x=N^{-1/6}$, and fix an
admissible choice of its auxiliary exponents and reference parameters.
For its direct, unmasked, and unseeded reference paths,
\cite[Proposition~6.1, (G1)]{OpenAILimitingLaw} states that there
exists $c_G>0$ such that
\begin{equation}\label{eq:intro-openai-G1}
 0\le F_*(\lambda)-f_N(b_\lambda,h_{*,\lambda})\le x^{4+c_G}
\end{equation}
for all sufficiently large $N$. Here $F_*$ is a scalar trial objective,
$f_N$ is a self-corrected hierarchical pressure evaluated at the clocks
$(b_\lambda,h_{*,\lambda})$, and $\lambda$ is a perturbation strength
that is unrelated to our tilt. We take terminal mass one and the
allowed strength $\lambda=x^{12}$. From the definitions
in~\cite[Section~3.1]{OpenAILimitingLaw}, one checks that the scalar
objective at $\lambda=0$ equals $P_\beta-\log2-T/4$. Using the perturbation estimates for the direct
scalar lift~\cite[Section~3.2, proof of Lemma~3.2]{OpenAILimitingLaw}
and the construction of the clocks~\cite[Section~3.4,
(P3)--(P4)]{OpenAILimitingLaw}, we find
\[
 \|b_\lambda-T\|_\infty+\|h_{*,\lambda}\|_\infty\le C_\beta\lambda,
 \qquad
 \Bigl|F_*(\lambda)-\Bigl(P_\beta-\log2-\frac T4\Bigr)\Bigr|
 \le C_\beta\lambda.
\]
With the uniform spin prior and centered Gaussian completion
of~\cite[Section~2.1]{OpenAILimitingLaw}, the unperturbed clocks satisfy
$f_N(T,0)=p_{N,\beta}-\log2-T/4$. Integrating the pressure variation
formula~\cite[Lemma~2.1, (A1)]{OpenAILimitingLaw} along the segment
between the two pairs of clocks, and using that its integrands lie in
$[0,1]$ at terminal mass one, we obtain
\[
 \bigl|f_N(b_\lambda,h_{*,\lambda})-f_N(T,0)\bigr|
 \le\frac14\|b_\lambda-T\|_1+\frac12\|h_{*,\lambda}\|_1
 \le C_\beta\lambda.
\]
We write $P_\beta-p_{N,\beta}$ as the sum of
$P_\beta-\log2-T/4-F_*(\lambda)$,
$F_*(\lambda)-f_N(b_\lambda,h_{*,\lambda})$, and
$f_N(b_\lambda,h_{*,\lambda})-f_N(T,0)$. We bound the first and last
differences by $C_\beta\lambda$ and the middle one by
\eqref{eq:intro-openai-G1}. Since $x^{4+c_G}=N^{-2/3-c_G/6}$ and
$x^{12}=N^{-2}$, we obtain
\begin{equation}\label{eq:intro-openai-bias}
 0\le P_\beta-p_{N,\beta}
 \le x^{4+c_G}+C_\beta x^{12}
 \le C_\beta N^{-2/3-\delta_\beta},\qquad
 \delta_\beta=c_G/6>0,
\end{equation}
after decreasing $c_G$ so that $c_G<8$ if necessary. The lower bound is
Guerra's, and \eqref{eq:intro-openai-bias} holds for the finitely many
remaining $N\ge2$ after we enlarge $C_\beta$. This bound follows from
\eqref{eq:intro-openai-G1} and the perturbation estimates above; it is
not a theorem on the mean bias stated in either paper. The proof
of~\cite[Proposition~6.1]{OpenAILimitingLaw} leaves $c_G$ unspecified,
and we do not extract from it the saving $c_G\ge1/5$ that the exponent
$7/10$ would require. The budgets of order $x^{5+c}$
in~\cite[Proposition~4.1]{OpenAILimitingLaw} concern increments of a
centered pressure~\cite[(A4)]{OpenAILimitingLaw}, from which the
finite-size mean has been subtracted, so they do not bound the mean
bias.

\paragraph{Conditional lower bounds.}
In the other direction, the variance--mean-deficit inequality
of~\cite[Lemma~7.4 and Remark~7.14, (7.68)]{AL} states that for every
$N\ge1$,
\[
 P_\beta-p_{N,\beta}\ge c_\beta
 \Bigl(\frac{\Var(F_{N,\beta})+\beta^2/2}{N}\Bigr)^{5/4}.
\]
Inserting $\Var(F_{N,\beta})\ge N^{1/3-4\varepsilon/5}$, we see that the
variance exponent stated in~\cite[Theorem~1.1]{OpenAIFluctuationScale}
would imply, for every fixed $\varepsilon>0$, that
$P_\beta-p_{N,\beta}\ge c_{\beta,\varepsilon}N^{-5/6-\varepsilon}$ for
all sufficiently large $N$. The positive limit of
$N^{-1/3}\Var(F_{N,\beta})$ stated
in~\cite[Theorem~1.2]{OpenAILimitingLaw} would give the stronger bound
\begin{equation}\label{eq:intro-openai-lower}
 P_\beta-p_{N,\beta}\ge c_\beta N^{-5/6}
 \qquad\text{for all sufficiently large }N.
\end{equation}
For \eqref{eq:intro-openai-lower}, an eventual lower bound
$\Var(F_{N,\beta})\ge k_\beta N^{1/3}$ with $k_\beta>0$ would suffice.
Combined with \cref{thm:main}, the variance exponent
of~\cite{OpenAIFluctuationScale} would imply that if
$P_\beta-p_{N,\beta}$ is of exact order $N^{-a}$, then $7/10\le a\le5/6$.

\paragraph{Verification status.}
We have checked the statements quoted above and the conditional
deductions from them. We have not independently reconstructed complete
proofs of the global comparison and fluctuation estimates
of~\cite{OpenAIFluctuationScale,OpenAILimitingLaw}, and we do not
assume their conclusions. The repository also contains Lean sources on
the fluctuations of the SK model. Their main theorem on the fluctuation
exponents, as we inspected it, takes as an explicit hypothesis a
two-sided bound on the centered logarithmic Laplace transform of the
completed free energy, \texttt{PositiveLaplaceWindow}; in the
fluctuation-scale paper, this bound is the conclusion
of~\cite[Proposition~7.5]{OpenAIFluctuationScale}. The conclusion of
that theorem concerns the fluctuation exponents, not a positive variance
prefactor or a limiting law.\footnote{We inspected the 151 files of the
directory \nolinkurl{lean/OAI/Probability/SKFluctuations} at commit
\nolinkurl{adc7f1241b42e322a6451854ab7e4b4c146bf78a}. The theorem
\nolinkurl{main_of_positiveLaplaceWindow}, at line~294 of
\nolinkurl{EventuallyNatPos.lean}, takes a hypothesis
\texttt{hw} of type \nolinkurl{PositiveLaplaceWindow}, which is defined at
lines~453--459 of \nolinkurl{Spin.lean}; the conclusion of the
theorem, \nolinkurl{MainFluctuationStatement}, is defined at
lines~531--540 of the same file. We found no proof supplying the
hypothesis in this directory. The directory was unchanged at commit
\nolinkurl{fd4aeeb2ee4fc729c18d98444fed42fd0529eeeb}, which we checked on
October~9, 2026. This was an inspection of the sources, not a
compilation.} Neither this conditional formal statement nor our
incomplete reconstruction shows that the stated results are false.
Our inspection did not identify an invalid step in the portions
examined.

\paragraph{Proof techniques.}
The limiting-law paper contains a one-site cavity expansion whose
structure is similar to that of \cref{part:cavity}.
In~\cite[Lemma~2.2]{OpenAILimitingLaw}, one site is chosen by
exchangeability and removed, and its Gaussian kernels are interpolated
to those of an independent scalar model on one site; the linear cavity
kernel is half of a connected four-spin correlation of the removed
spin~\cite[(A8)]{OpenAILimitingLaw}. In the proof
of~\cite[Proposition~6.1]{OpenAILimitingLaw}, the resulting linear
equation is inverted under a gap condition on the scalar
response~\cite[Proposition~5.1, (H1)]{OpenAILimitingLaw},
and~\cite[Lemma~5.2, (H2)]{OpenAILimitingLaw} writes the response on
the additive sector as a susceptibility multiplier minus a positive
semidefinite Hessian kernel. In our argument, the cavity inequality
contracts because the scalar response operator is at most $1+\eps_N$
with $\eps_N\le\rho/4$, while the interpolation contributes the factor
$1-\rho$. We derive the bound on the response operator from the second
variation of the Parisi functional under transports of levels
(\cref{lem:Jdom}); this operator acts on a space of transports, not on
the additive sector of~\cite{OpenAILimitingLaw}. The fluctuation-scale
paper also contains a one-site closure for the mean
overlap~\cite[Proposition~4.1]{OpenAIFluctuationScale}. We use two
tools from that paper. \Cref{part:cavity} relies on normalized replica
trees and covariance insertions of the type developed
in~\cite[Section~2]{OpenAIFluctuationScale} and used
in~\cite[Sections~2.1--2.3]{OpenAILimitingLaw}. In
\cref{lem:fnd-posterior}, we adapt the joint-source Brascamp--Lieb
argument of the proof of~\cite[Lemma~2.3]{OpenAIFluctuationScale}, which
is also used in the proof of~\cite[Lemma~2.1]{OpenAILimitingLaw}, to
squared conditional means on selected coordinates and to anisotropic
laws.



\section{Organization and conventions}\label{sec:intro-organization}

\Cref{sec:overview} contains an overview of the proof and a map of its
dependencies. The proof itself occupies five parts.
\begin{itemize}
\item \Cref{part:interpolation} concerns the interpolation.
\Cref{sec:fnd} recalls the interpolation of~\cite{AL}, its tilted law,
and the decomposition of the defect, and reduces the mean bias to
the fractional-moment error. \Cref{sec:gauss} proves two Gaussian tools
for tilted laws defined by arbitrary profiles, which we use in
\cref{part:profile,part:cavity}.
\item \Cref{part:overlap} contains overlap estimates at a minimizing
profile. \Cref{sec:wi} estimates moments of the self-overlap on
quantile intervals, and \cref{sec:th} bounds the mean--covariance term.
\item \Cref{part:profile} concerns the minimizing profile and the scalar
response. \Cref{sec:lz} shows that the minimizing profile is close in
$L^1$ to the scalar trial profile, and \cref{sec:lz-root} localizes its
lowest quantiles. \Cref{sec:sc} studies the scalar response operator and
proves its stability at the minimizing profile.
\item \Cref{part:cavity} contains the cavity argument. \Cref{sec:c7}
sets up the cavity interpolation and its tree laws,
\cref{sec:c7-endpoint} computes the derivative at the independent
endpoint, and \cref{sec:c7-closure} closes the cavity inequality and
bounds the defect.
\item In \cref{part:proof}, \cref{sec:b7} integrates these estimates and
proves \cref{thm:main}.
\end{itemize}

We rely on~\cite{AL} for the interpolation identities and the scalar
floor estimates, and on~\cite{Lopatto} and~\cite{AuffingerChen2015Properties}
for the structure and regularity of $\mu_\beta$. We restate these
results where they are used. Locators in~\cite{AL} refer to its arXiv
version~1, locators in~\cite{Lopatto} to its arXiv version~4, and
locators in the OpenAI manuscripts to their initial repository version.

Throughout, $\beta>1$ is fixed, and $C_\beta$ denotes a constant
depending only on $\beta$ whose value may change from line to line; a
numerical constant does not depend on any parameter. All profiles defining a tilted law are
deterministic: minimizers may depend on $N$, $\rho$, and $\lambda$ but
not on the couplings. For symmetric matrices or self-adjoint operators,
$B\preceq C$ means that $C-B$ is positive semidefinite. We write
$\log_+x=\log(\max\{1,x\})$ for $x\ge0$. Distribution functions, such
as $\alpha$ and the profiles $a$ of \cref{sec:fnd}, are
right-continuous. Atoms of measures are retained in Stieltjes
integrals, and an interval $(q,1]$ excludes an atom at $q$.
